\chapter{Modelo de rendimiento y notación}

La utilidad de una arquitectura de datos no puede deducirse únicamente de la claridad
de sus clases o diagramas. Una separación correcta de responsabilidades puede reducir
errores y, al mismo tiempo, introducir lecturas, escrituras, almacenamiento y tareas de
mantenimiento adicionales. Esta parte establece un modelo para razonar sobre esos
compromisos antes de medirlos y un protocolo para comprobarlos sin atribuir a COBRIA
resultados que dependan del motor, la red o una configuración particular.

El análisis distingue dos preguntas. La primera es asintótica: cómo crece el trabajo
cuando aumenta el volumen de datos. La segunda es empírica: cuánto tarda y cuánto cuesta
una operación bajo una carga y una plataforma concretas. La notación $O(\cdot)$ responde
a la primera pregunta; no representa milisegundos, precio monetario ni garantía de
rendimiento. Dos operaciones de orden $O(1)$ pueden tener costos reales muy distintos
si una requiere una lectura local y otra varias comunicaciones de red
\parencite{cormen2022algorithms}.

\section{Variables del modelo}

Sea $N$ el número de entidades canónicas dentro de un ámbito; $R$ el número de resultados
de una consulta; $P$ el número de proyecciones asociadas con una entidad; $P_w\leq P$ el
número de proyecciones afectadas por una escritura; $\Delta N$ el número de entidades
modificadas desde un marca de corte; $B_c$ el tamaño medio de un canónico y $B_{p_i}$ el
tamaño medio de su contribución a la proyección $i$. Se utilizan además:

\begin{table}[H]
\centering
\caption{Símbolos del modelo de rendimiento.}
\label{tab:simbolos-rendimiento}
\begin{tabularx}{\textwidth}{>{\bfseries}p{2cm}X>{\bfseries}p{2cm}X}
\toprule
Símbolo & Significado & Símbolo & Significado \\
\midrule
$L$ & latencia observada & $Q$ & operaciones por unidad de tiempo \\
$S$ & bytes almacenados & $T$ & bytes transferidos \\
$W_f$ & escrituras físicas & $R_f$ & lecturas físicas \\
$C_x$ & costo unitario del recurso $x$ & $h$ & proporción de aciertos de caché \\
$f_q$ & frecuencia de la consulta & $f_w$ & frecuencia de escritura \\
$d$ & número de dependencias & $k$ & tamaño del lote \\
\bottomrule
\end{tabularx}
\end{table}

Estas variables se calculan por organización y también de manera agregada. El promedio global
puede ocultar que un organización pequeño recibe peor servicio o que uno grande domina el
consumo. Por ello, cada resultado debe conservar al menos identificador anonimizado de
ámbito, tipo de operación, versión del contrato y configuración experimental.

\section{Frontera de medición}

Una medición debe indicar dónde comienza y termina. COBRIA propone tres fronteras:

\begin{enumerate}[leftmargin=*]
  \item \textbf{Algorítmica:} trabajo efectuado por la transformación, sin red ni
  persistencia.
  \item \textbf{Componente:} Repositorio, Gestor de proyecciones o Constructor de conjuntos de datos junto
  con sus dependencias simuladas o reales.
  \item \textbf{Extremo a extremo:} desde el caso de uso hasta la respuesta observable,
  incluyendo autorización, serialización, red, almacenamiento y reintentos.
\end{enumerate}

No deben compararse cifras de fronteras distintas como si midieran el mismo fenómeno.
La instrumentación se incorpora fuera del dominio siempre que sea posible para evitar
que el modelo canónico dependa de un proveedor de telemetría.

\section{Métricas principales}

La latencia se reportará mediante mediana y percentiles p95 y p99, no solo con el
promedio. El rendimiento efectivo se expresa como operaciones completadas por segundo bajo una
concurrencia declarada. También se registrarán bytes leídos y escritos, cantidad de
operaciones físicas, errores, reintentos, uso de memoria, tiempo de CPU y retraso de
actualización de proyecciones. El enfoque sigue la práctica de separar carga de trabajo,
métricas y configuración al comparar sistemas de almacenamiento
\parencite{cooper2010ycsb,jain1991performance}.

\chapter{Complejidad computacional de las operaciones}

\section{Acceso canónico}

Cuando una clave física se deriva directamente de $(S,id)$ y el motor ofrece acceso
indexado por clave, la búsqueda lógica de una entidad se modela como $O(1)$ promedio.
Esta expresión presupone una clave conocida; no cubre resolución DNS, autenticación,
red, deserialización ni contención. Si se busca por un atributo no indexado, el costo
puede crecer hasta $O(N)$ porque deben examinarse los registros del ámbito.

La recuperación de una página de $R$ resultados desde una proyección ordenada puede
modelarse como $O(\log N+R)$ cuando existe una estructura de búsqueda balanceada, o como
$O(R)$ después de localizar un cursor. El uso de offset puede degradarse al aumentar la
profundidad; COBRIA prefiere cursores estables compuestos por orden, ámbito e identidad.

\section{Creación y actualización}

Validar un objeto con $m$ campos tiene costo $O(m)$ si cada regla es local. Las reglas
que consultan otras entidades agregan $d$ accesos y deben expresarse como
$O(m+d)$ más el costo de esos accesos. Persistir solamente el canónico requiere una
escritura lógica. Mantener síncronamente $P_w$ proyecciones agrega, en el caso simple,
$O(P_w)$ transformaciones y escrituras:

\begin{equation}
T_{write}=T_{validate}+T_{canonical}+\sum_{i=1}^{P_w}
  (T_{g_i}+T_{persist_i}).
\label{eq:tiempo-escritura}
\end{equation}

La suma no implica que todos los términos sean iguales. Una proyección de clave puede
ser constante, mientras un agregado que inspecciona una colección puede depender de su
tamaño. La ejecución paralela reduce tiempo de pared, pero no elimina trabajo total ni
garantiza atomicidad.

\section{Reconstrucción total e incremental}

Una reconstrucción total que aplica $P$ transformaciones independientes a $N$ entidades
tiene costo superior $O(NP)$. Si cada entidad alimenta solamente un subconjunto fijo,
el costo efectivo es $O(N\bar{P})$, donde $\bar{P}$ es el promedio de proyecciones
aplicables. La memoria puede mantenerse en $O(k)$ mediante procesamiento por lotes, en
vez de cargar $O(N)$ objetos.

Una reconstrucción incremental basada en un índice confiable de cambios procesa
$\Delta N$ entidades:

\begin{equation}
T_{incremental}\approx O(\Delta N\bar{P})+T_{discover}+T_{verify}.
\end{equation}

Es ventajosa cuando $\Delta N\ll N$ y el costo de descubrir cambios no domina. Si el
marca de corte es incompleto, una reconstrucción incremental rápida puede producir un estado
incorrecto. La optimización, por tanto, está subordinada a la cobertura del registro de
cambios.

\section{Verificación de proyecciones}

Comparar cada canónico con sus derivados cuesta $O(N\bar{P})$. Un muestreo de $s$
entidades reduce el trabajo a $O(s\bar{P})$, pero solo estima la tasa de anomalías. Los
resúmenes criptográficos por partición permiten localizar diferencias sin transferir todos los objetos;
su construcción inicial sigue siendo lineal y su actualización requiere disciplina de
versionado.

\begin{table}[H]
\centering
\caption{Complejidad esperada de operaciones COBRIA.}
\label{tab:complejidad-operaciones}
\begin{tabularx}{\textwidth}{p{4cm}p{3.2cm}X}
\toprule
Operación & Orden esperado & Condición principal \\
\midrule
Buscar canónico por clave & $O(1)$ promedio & Ruta derivable de ámbito e ID.\\
Escanear por atributo & $O(N)$ & No existe proyección o índice.\\
Consultar proyección & $O(\log N+R)$ & Índice ordenado; depende del motor.\\
Validar entidad & $O(m+d)$ & Reglas locales y dependencias explícitas.\\
Actualizar derivados & $O(P_w)$ & Transformaciones acotadas por proyección.\\
Reconstrucción total & $O(N\bar{P})$ & Recorrido completo por lotes.\\
Reconstrucción incremental & $O(\Delta N\bar{P})$ & Registro de cambios completo.\\
Verificación completa & $O(N\bar{P})$ & Comparación fuente--derivado.\\
Generar conjunto de datos & $O(NF)$ & $F$ características por entidad, sin uniones externos.\\
\bottomrule
\end{tabularx}
\end{table}

La tabla representa modelos de crecimiento y no contratos del proveedor. Cada
implementación deberá confirmar sus supuestos con planes de consulta, contadores y
pruebas de carga.

\chapter{Amplificación y modelo de costos}

\section{Amplificación de lectura}

La amplificación de lectura relaciona trabajo físico y resultado útil. Para una consulta
$q$ se define:

\begin{equation}
A_r(q)=\frac{B_{read}(q)}{\max(B_{result}(q),\epsilon)},
\end{equation}

donde $\epsilon$ evita división por cero en consultas sin resultados. También puede
calcularse por registros: $A_r^{reg}=R_f/\max(R,1)$. Una consulta que escanea diez mil
registros para devolver diez tiene amplificación aproximada de mil. Una proyección bien
diseñada acerca el numerador al volumen útil, aunque añade mantenimiento y almacenamiento.

\section{Amplificación de escritura}

Para una operación lógica $o$:

\begin{equation}
A_w(o)=\frac{W_f(o)}{W_l(o)},
\end{equation}

donde normalmente $W_l=1$. Si se persiste el canónico, una entrada cola de salida y tres
proyecciones, $A_w=5$ sin contar réplicas internas del motor. Deben reportarse por
separado la amplificación visible para la aplicación y la amplificación interna cuando
el proveedor la exponga.

\section{Amplificación de almacenamiento}

El factor de almacenamiento es:

\begin{equation}
A_s=\frac{S_{canonical}+S_{projections}+S_{metadata}+S_{registros}}
{S_{canonical}}.
\end{equation}

Una proyección no se justifica por ser pequeña de forma aislada. Su costo acumulado
incluye versiones durante reconstrucciones, índices del motor, historial, respaldos y
transferencia. COBRIA asigna a cada proyección un presupuesto y un propietario.

\section{Costo total por ventana}

Para una ventana temporal $\tau$, el costo observable se modela como:

\begin{align}
C_{total}(\tau)=&\ C_{read}R_f + C_{write}W_f + C_{store}S\tau \\
&+ C_{transfer}T + C_{compute}U + C_{ops}H,
\label{eq:costo-total}
\end{align}

donde $U$ representa consumo computacional y $H$ esfuerzo operativo estimado. El último
término evita declarar óptima una solución barata para el proveedor pero costosa de
operar, depurar o reconstruir. Los precios cambian; por ello el experimento publicará
cantidades físicas y dejará los precios como parámetros reemplazables.

\section{Punto de equilibrio de una proyección}

Sea $C_s$ el costo de una consulta sin proyección, $C_p$ el costo de consultarla con
proyección y $C_m$ el costo de mantenerla por escritura. La proyección produce beneficio
en una ventana cuando:

\begin{equation}
f_q(C_s-C_p)>f_wC_m+C_{storage}+C_{rebuild}.
\label{eq:break-even}
\end{equation}

Esta desigualdad no pretende reducir toda decisión a dinero. Latencia máxima,
disponibilidad, aislamiento y corrección pueden actuar como restricciones obligatorias.
Sí permite identificar proyecciones con uso insuficiente o mantenimiento excesivo.

\section{Presupuesto de proyecciones}

Cada contrato debe declarar límites de bytes, escrituras por cambio, retraso tolerable,
tiempo máximo de reconstrucción y frecuencia mínima que justifica su existencia. Un
reporte periódico compara presupuesto y observación. Si una proyección excede el límite,
las opciones son optimizarla, reducir campos, agrupar eventos, cambiar consistencia o
retirarla; no se oculta el exceso mediante promedios.

\chapter{Estrategias de optimización}

\section{Optimizar desde la consulta}

La secuencia recomendada parte de una consulta observable: frecuencia, selectividad,
orden, tamaño de respuesta y objetivo de latencia. Después se decide si basta una clave
canónica, un índice del motor, una proyección específica o una caché. Crear estructuras
derivadas antes de conocer la carga aumenta $P$ y puede empeorar H9.

Una proyección debe almacenar solo campos necesarios para filtrar, ordenar y responder.
Duplicar el documento completo facilita algunas lecturas, pero aumenta transferencia,
amplificación y riesgo de exposición. La paginación por cursor evita recorrer desplazamientos
crecientes y conserva una frontera consistente de lectura cuando el orden incluye una
clave de desempate.

\section{Lotes y paralelismo controlado}

El tamaño de lote $k$ equilibra memoria, viajes de red y duración de transacciones. Un
lote demasiado pequeño multiplica comunicaciones; uno demasiado grande presiona memoria
y extiende la recuperación después de un fallo. COBRIA no fija un $k$ universal: lo
selecciona experimentalmente y registra la configuración.

El paralelismo se limita por partición, cuota y capacidad de escritura. Aumentar procesos de trabajo
puede reducir tiempo hasta alcanzar contención o limitación de capacidad; después, el rendimiento efectivo se
estanca y la latencia de cola aumenta. El controlador debe aplicar control de contrapresión, variación aleatoria
y reintentos acotados, conservando idempotencia.

\section{Caché con ámbito y revisión}

Con una tasa de aciertos $h$, el costo medio simplificado de lectura es:

\begin{equation}
E[T]=hT_{hit}+(1-h)T_{miss}.
\end{equation}

La clave incluye organización, tipo, ID y, cuando corresponda, versión. La caché no sustituye
una proyección si se necesita consulta por atributos, ni se considera autoridad. TTL
reduce permanencia, pero no garantiza invalidación inmediata. Para datos sensibles se
deben medir además exposición, borrado y segregación.

\section{Reconstrucción eficiente}

Una reconstrucción segura y optimizada sigue cinco decisiones:

\begin{enumerate}[leftmargin=*]
  \item fijar versión y marca de corte de entrada;
  \item dividir por ámbito y rango estable;
  \item procesar lotes idempotentes con checkpoints;
  \item escribir en una versión paralela de la proyección;
  \item verificar y efectuar un cambio de versión atómico o reversible.
\end{enumerate}

La versión paralela requiere almacenamiento temporal, pero evita exponer una mezcla de
contratos. Para reconstrucciones incrementales se conserva una operación total como
mecanismo de recuperación y como prueba periódica de reconstruibilidad.

\section{Optimización adaptativa}

La arquitectura puede recomendar cambios a partir de telemetría, sin permitir que una
heurística modifique por sí sola el estado canónico. Un ciclo adaptativo observa
frecuencia, amplificación, latencia y retraso; propone crear, rediseñar o retirar una
proyección; estima el punto de equilibrio; ejecuta una comparación controlada; y somete
el cambio a aprobación. Esto hace posible usar minería de datos para descubrir patrones
de acceso y aprendizaje automático para predecir carga, manteniendo una frontera de
gobierno.

\begin{figure}[H]
\centering
\begin{tikzpicture}[
  node distance=1.15cm and 1.25cm,
  box/.style={draw,rounded corners,align=center,minimum width=2.55cm,
    minimum height=.9cm,fill=cobriaLight},
  line/.style={-{Latex[length=2mm]},thick}
]
\node[box] (observe) {Observar\\telemetría};
\node[box,right=of observe] (analyze) {Analizar\\costo y carga};
\node[box,right=of analyze] (propose) {Proponer\\cambio};
\node[box,below=of propose] (evaluate) {Evaluar\\en paralelo};
\node[box,left=of evaluate] (approve) {Aprobar y\\versionar};
\node[box,left=of approve] (verify) {Verificar\\resultado};
\draw[line] (observe) -- (analyze);
\draw[line] (analyze) -- (propose);
\draw[line] (propose) -- (evaluate);
\draw[line] (evaluate) -- (approve);
\draw[line] (approve) -- (verify);
\draw[line] (verify) -- (observe);
\end{tikzpicture}
\caption{Ciclo gobernado de optimización adaptativa.}
\label{fig:ciclo-optimizacion}
\end{figure}

\chapter{Protocolo experimental de rendimiento}

\section{Objetivo y comparaciones}

El protocolo evaluará si COBRIA reduce el costo de lectura y mejora reconstrucción y
trazabilidad a cambio de costos cuantificables. No se presentarán resultados hasta
ejecutar el protocolo. Se compararán, como mínimo, estas configuraciones:

\begin{enumerate}[leftmargin=*]
  \item \textbf{B0, acceso directo:} canónicos sin proyección específica;
  \item \textbf{B1, optimización nativa:} índices disponibles en el motor;
  \item \textbf{C1, COBRIA síncrona:} canónico y proyecciones críticas en la operación;
  \item \textbf{C2, COBRIA asíncrona:} canónico, cola de salida y proyecciones actualizadas por
  procesos de trabajo;
  \item \textbf{C3, COBRIA incremental:} C2 con registro de cambios y reconstrucción
  incremental.
\end{enumerate}

No todas las configuraciones aplican a toda operación. La comparación se declara antes
de observar resultados y mantiene la misma semántica, conjunto de datos y política de consistencia.

\section{Factores y niveles}

\begin{table}[H]
\centering
\caption{Factores propuestos para los experimentos.}
\label{tab:factores-experimento}
\begin{tabularx}{\textwidth}{p{3.6cm}X X}
\toprule
Factor & Niveles iniciales & Propósito \\
\midrule
Volumen $N$ & $10^3$, $10^4$, $10^5$ y, si es viable, $10^6$ & Estimar tendencia de
crecimiento.\\
Mezcla lectura/escritura & 95/5, 80/20, 50/50 & Representar cargas distintas.\\
Concurrencia & 1, 8, 32 y 128 clientes & Detectar saturación y contención.\\
Número de proyecciones & 0, 1, 3, 5 y 10 & Medir amplificación marginal.\\
Distribución & uniforme y sesgada & Evaluar claves calientes.\\
Tamaño de entidad & pequeño, medio y grande & Medir serialización y transferencia.\\
Cambios $\Delta N/N$ & 0.1\%, 1\%, 10\% y 100\% & Comparar reconstrucción incremental
y total.\\
\bottomrule
\end{tabularx}
\end{table}

Los valores son puntos iniciales y se ajustarán a la capacidad del entorno antes de
registrar la especificación definitiva. El Sistema A aportará cargas temporales,
multimodales y de eventos; el Sistema B, operaciones empresariales, catálogos e
aislamiento multiinquilino. Los nombres y datos originales no aparecerán en scripts ni
resultados.

\section{Datos y cargas reproducibles}

Los generadores usarán semilla publicada, distribuciones declaradas y relaciones
válidas. Para evitar que datos triviales favorezcan una estrategia, incluirán cardinalidad,
sesgo, valores ausentes, revisiones y tamaños variables. Los artefactos multimedia se
representarán mediante metadatos y archivos sintéticos; no se copiará contenido privado.

Cada ejecución incluye calentamiento, ventana de medición y enfriamiento. Se aleatoriza
el orden de configuraciones cuando sea posible y se repiten ejecuciones independientes.
Se registra hardware, sistema operativo, entorno de ejecución, versión del motor, región, límites,
configuración de índices y revisión del código. Un prueba comparativa como YCSB sirve de referencia
para estructurar cargas, pero las operaciones de dominio y reconstrucción requieren
escenarios COBRIA adicionales \parencite{cooper2010ycsb}.

\section{Instrumentación y control}

Cada operación recibe un identificador de correlación. Los contadores se toman en el
cliente y, cuando exista acceso, en el almacenamiento. Se sincronizan relojes para medir
retraso asíncrono, pero la latencia local utiliza un reloj monotónico. Se evita registrar
cargas útiles sensibles. La telemetría debe distinguir éxito, rechazo de negocio, conflicto,
timeout y reintento.

Los ensayos controlarán cachés, calentamiento, concurrencia y tareas de fondo. Si un
servicio administrado introduce ruido no controlable, se reportará y se aumentarán
repeticiones, sin eliminar observaciones solo por ser desfavorables.

\section{Análisis estadístico}

Para cada combinación se reportarán tamaño muestral, mediana, p95, p99, intervalo de
confianza y distribución gráfica. Las comparaciones incluirán diferencia absoluta,
razón y tamaño de efecto. Antes de usar pruebas paramétricas se examinarán independencia,
forma de distribución y varianza. Cuando no se satisfagan supuestos, se preferirán
intervalos remuestreo con reemplazo o pruebas no paramétricas. Las múltiples comparaciones se
identificarán y corregirán cuando corresponda.

La significancia estadística no reemplaza relevancia práctica. Una mejora pequeña puede
carecer de valor si aumenta mucho $A_w$, $A_s$ o el retraso. El análisis de sensibilidad
variará precios, frecuencias y objetivos de latencia en la ecuación
\ref{eq:costo-total}.

\section{Criterios de aceptación}

Las hipótesis se juzgarán con reglas previas:

\begin{itemize}[leftmargin=*]
  \item H1 se apoya si el crecimiento de lecturas se aproxima a $R$ y la latencia de
  cola mejora sin cambiar la semántica;
  \item H2 se cuantifica, no se considera un fallo: toda reducción se acompaña de
  $A_w$ y $A_s$;
  \item H3 requiere igualdad según el contrato, ausencia de huérfanos y cobertura
  completa después de reconstruir;
  \item H4 compara total e incremental para distintos $\Delta N/N$ y busca su punto de
  cruce;
  \item H9 se apoya si existen regiones donde el costo marginal supera el beneficio de
  lectura.
\end{itemize}

No se modificará un umbral después de conocer resultados sin declarar el cambio como
exploratorio.

\section{Plantilla de resultados}

\begin{table}[H]
\centering
\caption{Matriz mínima de resultados por escenario.}
\label{tab:plantilla-resultados}
\begin{tabularx}{\textwidth}{p{3.7cm}X X}
\toprule
Dimensión & Indicadores & Interpretación requerida \\
\midrule
Respuesta & p50, p95, p99, rendimiento efectivo & Saturación y variabilidad.\\
Lectura & registros, bytes, $A_r$ & Trabajo evitado frente a resultado.\\
Escritura & operaciones, bytes, $A_w$ & Precio de mantener derivados.\\
Almacenamiento & bytes y $A_s$ & Incluye versiones e índices.\\
Reconstrucción & duración, cobertura, anomalías & Exactitud antes que velocidad.\\
Consistencia & retraso, conflictos, reintentos & Ventana de obsolescencia.\\
Operación & CPU, memoria, horas de trabajo & Costo total y mantenibilidad.\\
Equidad de ámbito & percentiles por organización & Evitar ocultar degradación.\\
\bottomrule
\end{tabularx}
\end{table}

\chapter{Rendimiento para minería de datos e inteligencia artificial}

\section{Preparación analítica como proyección}

Un conjunto de datos se trata como una proyección versionada del estado canónico y de eventos
permitidos. Su costo de construcción depende de entidades, ventanas temporales y número
de características. Para $N$ entidades y $F$ características locales, el caso simple es
$O(NF)$. Uniones, agregados temporales, texto o embeddings añaden términos que deben
medirse por separado.

Esta estructura favorece minería de datos porque conserva el significado de la entidad,
el ámbito, el tiempo del evento, la versión de transformación y el linaje. Permite volver
a producir un conjunto de datos, comparar versiones y detectar qué cambios de datos explican una
diferencia. No garantiza que las características sean útiles ni que el modelo sea justo.

\section{Corrección temporal y prevención de fuga}

Para una observación con tiempo de corte $t$, solo se utilizan hechos disponibles en o
antes de $t$. Una característica calculada con datos futuros produce fuga y métricas de
entrenamiento engañosas. El Constructor de conjuntos de datos recibe marca de corte, ventana y versión; el
resultado registra tiempo de evento, tiempo de ingestión y tiempo de cálculo.

El acceso punto-en-tiempo puede requerir historial canónico, log de eventos o snapshots.
Si el sistema conserva únicamente el estado actual, COBRIA debe declarar que ciertas
reconstrucciones históricas no son posibles en vez de inferirlas.

\section{Características, embeddings y multimodalidad}

Las características tabulares, series, texto y referencias multimedia pueden compartir
un contrato de procedencia. Los archivos grandes permanecen fuera del objeto canónico;
este conserva ubicación autorizada, hash, tipo, versión y política. Una extracción de
audio, texto o imagen genera un derivado con modelo, parámetros y fecha.

Los embeddings son proyecciones reconstruibles si se conserva la entrada autorizada,
el identificador exacto del modelo y la transformación. Cambiar de modelo produce una
nueva versión; no se mezclan vectores incompatibles en el mismo índice sin una estrategia
explícita. La reconstrucción puede ser costosa, por lo que se presupuestan cómputo,
transferencia y ventana de coexistencia.

\section{RAG y agentes}

En recuperación aumentada, el índice vectorial no es autoridad: sirve para localizar
candidatos. La respuesta debe enlazar fragmento, documento canónico, versión y permisos.
La autorización se aplica antes de incorporar contexto al instrucción. Un agente puede leer
proyecciones adaptadas a su tarea, pero cualquier cambio de negocio vuelve a un caso de
uso que valida identidad, ámbito, invariantes e idempotencia.

Esta frontera limita el daño de una inferencia incorrecta y permite auditar propuesta,
evidencia, decisión humana y efecto final. El rendimiento del modelo se mide separado
del rendimiento de datos: recuperación, construcción de contexto, inferencia y
persistencia tienen latencias y costos distintos.

\section{Optimización de cadenas de procesamiento de IA}

COBRIA permite reutilizar transformaciones deterministas, materializar características
de alto uso y reconstruirlas en lotes. La decisión se rige por el mismo punto de
equilibrio de la ecuación \ref{eq:break-even}. Para entrenamiento se favorecen lotes
grandes y reproducibilidad; para inferencia en línea, baja latencia y frescura. Una sola
representación física rara vez optimiza ambos caminos, por lo que los contratos deben
declarar equivalencia semántica entre versión fuera de línea y en línea, idea central de los
Almacenes de características \parencite{orr2021featurestores,li2023featurestore}.

\begin{figure}[H]
\centering
\begin{tikzpicture}[
  node distance=.95cm and 1.05cm,
  box/.style={draw,align=center,minimum width=2.45cm,minimum height=.9cm},
  derived/.style={box,fill=cobriaLight},
  line/.style={-{Latex[length=2mm]},thick}
]
\node[box] (canonical) {Datos canónicos\\y eventos};
\node[derived,right=of canonical] (builder) {Constructor de conjuntos de datos\\versionado};
\node[derived,right=of builder] (datosNodo) {Conjunto de datos /\\Conjunto de características};
\node[derived,below=of datosNodo] (model) {Modelo y\\evaluación};
\node[derived,left=of model] (prediction) {Predicción con\\evidencia};
\node[box,left=of prediction] (usecase) {Caso de uso\\autorizado};
\draw[line] (canonical) -- (builder);
\draw[line] (builder) -- (datosNodo);
\draw[line] (datosNodo) -- (model);
\draw[line] (model) -- (prediction);
\draw[line] (prediction) -- (usecase);
\draw[line,dashed] (usecase) -- node[left,font=\small]{cambio validado} (canonical);
\end{tikzpicture}
\caption{Flujo gobernado desde datos canónicos hasta decisiones asistidas por IA.}
\label{fig:flujo-ia}
\end{figure}

\section{Métricas analíticas y de IA}

Además de exactitud, precisión, exhaustividad u otras métricas propias del problema, se
registrarán cobertura de linaje, reproducibilidad entre ejecuciones, frescura,
divergencia en línea--fuera de línea, proporción de campos desconocidos, costo por ejemplo y
latencia por etapa. Para RAG se separan recuperación y generación; para agentes se miden
también acciones rechazadas, confirmaciones y efectos reversibles. La gestión de riesgo
acompaña el ciclo completo y no se reduce a una puntuación del modelo
\parencite{nist2023airmf}.

\chapter{Amenazas a la validez y síntesis de la Parte III}

\section{Validez interna}

Cambios de caché, calentamiento, throttling, tareas de fondo o distinta configuración
pueden explicar una diferencia atribuida a COBRIA. Se mitigarán mediante automatización,
orden aleatorio, repetición, configuración registrada y comparación sobre la misma
semántica. La instrumentación también tiene costo; se medirá una configuración de
control cuando su efecto pueda ser relevante.

\section{Validez de constructo}

Latencia no equivale a calidad arquitectónica, cantidad de clases no equivale a
mantenibilidad y reconstrucción exitosa no equivale a corrección del negocio. Por ello
cada concepto utiliza varios indicadores. La reconstruibilidad exige fuente suficiente,
contrato versionado, resultado verificable y ausencia de autoridad en el derivado.

\section{Validez externa}

Dos sistemas anonimizados permiten variación de dominio, pero no representan todos los
motores ni escalas. Los emuladores pueden diferir de servicios administrados. Los
resultados se generalizarán como condiciones y mecanismos ---por ejemplo, relación entre
frecuencia de lectura y amplificación---, no como una promesa universal de milisegundos.

\section{Validez de conclusión}

Muestras pequeñas, colas largas y múltiples comparaciones pueden producir conclusiones
inestables. Se publicarán distribuciones, intervalos, tamaños de efecto y resultados
negativos. Una hipótesis no apoyada no se ocultará; puede delimitar el dominio donde la
arquitectura deja de ser conveniente.

\section{Síntesis de la Parte III}

Esta parte convirtió la afirmación general de ``optimizar'' en magnitudes observables.
El beneficio esperado de las proyecciones es reducir el trabajo de lectura desde una
dependencia de $N$ hacia una dependencia más cercana a $R$, a cambio de trabajo de
escritura $O(P_w)$, almacenamiento adicional y reconstrucción. Las ecuaciones de
amplificación y costo permiten comparar ese intercambio sin confundir complejidad
asintótica con tiempo real.

También se definieron estrategias de optimización, un protocolo reproducible y reglas
de análisis que impiden inventar resultados. Finalmente, el mismo principio de
proyección reconstruible se extendió a conjuntos de datos, características, embeddings, RAG y
agentes, incorporando corrección temporal, linaje y autorización. Con ello queda
preparada la siguiente parte: desarrollar en profundidad la adaptabilidad de COBRIA a
analítica e inteligencia artificial y formalizar sus mecanismos de gobierno.
